\documentclass[12pt]{article}

\usepackage[margin=1in]{geometry}
\usepackage{setspace}
\usepackage{microtype}
\usepackage{parskip}
\usepackage[backend=biber,style=numeric,sorting=nyt]{biblatex}
\addbibresource{references.bib}

\onehalfspacing

\title{Washington University PhD Personal Statement}
\author{Rowan Torbitzky-Lane}
\date{December 1, 2026}

\begin{document}

\maketitle

I embarked on my first challenging computer project when I was ten. I wanted to start my own Minecraft server. After learning the steps from a YouTube video and forwarding a port, I was running a Minecraft server on my dad's old laptop. As I grew older, my curiosity broadened. I learned to program in high school and built my own PC. After attempts to study other fields, my place is in computer vision or graphics.

My professional experience ranges machine learning with relevant applications to computer vision/graphics. To begin, my first professional experience was at Hunter Engineering. I learned PyTorch, Anaconda, and used \textit{YOLOv5} \cite{terven2023yolo} to recognize license plates from North America and Europe. Furthermore, I wrote the first iteration of an automated data downloading process still in use today. The next year, I got my first taste of research. I worked at the Center for Intelligent Infrastructure at Missouri S\&T. I was tasked with using optical flow to propagate image segmentation annotations in a video. We used outdated MATLAB code which needed updating to run on the latest version. I had a basic prototype functional before I left. I did not see eye to eye with my coworker. Looking back, I would handle the situation differently and not hastily leave. I was accepted into a REU program the next summer at Missouri S\&T. I learned how to read papers, gained basic knowledge of transformers, and experienced the world of bioinformatics. I read \textit{Attention is All You Need} \cite{vaswani2017attention} and papers discussing bi-directional LLMs \cite{schiff2024caduceus}. My latest research experience was conducted at USC's Institute for Creative Technologies. I was tasked with giving gaussian splats physics properties in conjunction with an underlying mesh inside of the Unity game engine. We needed to efficiently correspond about 7 million splats (which are points in three dimensional space with extra metadata) to a mesh with about 350 thousand faces. C\# was too slow of a language (about 1000 spalts processed per second). Using a Rust FFI, I 

\pagebreak
\printbibliography

\end{document}
